\documentclass[conference]{IEEEtran}
\IEEEoverridecommandlockouts
\usepackage{cite}
\usepackage{amsmath,amssymb,amsfonts}
\usepackage{algorithmic}
\usepackage{algorithm}
\usepackage{graphicx}
\usepackage{subcaption}
\usepackage{textcomp}
\usepackage{xcolor}
\usepackage{booktabs}
\usepackage{multirow}
\usepackage{url}
\usepackage{hyperref}
\hypersetup{
  colorlinks=true,
  linkcolor=black,
  citecolor=black,
  urlcolor=blue
}

\def\BibTeX{{\rm B\kern-.05em{\sc i\kern-.025em b}\kern-.08em
    T\kern-.1667em\lower.7ex\hbox{E}\kern-.125emX}}

\begin{document}

\title{The Attacker in the Gap: ECMP Salt-Collision Link Flooding and a Rotation-Cadence Defence}

\author{\IEEEauthorblockN{Peter Ja\v{s}}
\IEEEauthorblockA{\textit{Department of Computers and Informatics} \\
\textit{Technical University of Ko\v{s}ice}\\
Ko\v{s}ice, Slovakia \\
peterjas81@gmail.com}
}

\maketitle

\begin{abstract}
Equal-Cost Multi-Path (ECMP) routing selects an egress link by hashing the packet 5-tuple together with a
controller-managed salt: $\mathrm{link} = h(\mathrm{5\text{-}tuple}, \mathrm{salt}) \bmod N$. When that salt is
generated by a weak, reconstructable source, a single attacker who recovers it can craft many individually
legitimate-looking flows that all resolve to one victim-carrying egress link, saturating it while every other
link in the fabric reports healthy. We show this attacker is not merely more precise than a volumetric flood --
it is structurally invisible to the two mechanisms that define current data-plane DoS practice. Per-source rate
limiting cannot catch it because no single source exceeds a bandwidth cap; the damage is done by \emph{placement},
not \emph{quantity}. Per-source connection throttling cannot catch it either, because the attacker deliberately
spreads its collision set across many distinct, individually-compliant 5-tuples. We implement and evaluate the
attack and its evasion of both defences by driving the real, frozen hash-core, salt-source, and defence-policy
implementations through a deterministic flow-level simulation of a controller-managed ECMP fabric -- a
preliminary study whose direct next step is a bare-metal implementation -- and evaluate a mechanism-level
defence -- rotating the hash salt at the controller -- that invalidates the attacker's pre-computed collision
set without inspecting any traffic. We further derive and confirm a
rotation-frequency specification that maps an attacker's salt-reconstruction time to the minimum safe rotation
interval, giving operators an actionable cadence rather than a qualitative recommendation. Swapping the salt
source from a cryptographically secure PRNG to a quantum random number generator (QRNG) served over a
provenance-attesting entropy-as-a-service endpoint produces an outcome statistically indistinguishable from the
CSPRNG case under this threat model -- an honest null result -- but the QRNG-as-a-service channel adds a
practical, auditable capability a CSPRNG cannot: a signed, per-draw entropy provenance receipt suited to
compliance- and attestation-sensitive multi-tenant deployments. We position the contribution carefully against
adjacent prior art -- the RSS secret-key defence for NIC receive-queue flooding, single-shard flooding in
sharded systems, and the Crossfire/Coremelt line of source-indistinguishable link-flooding attacks -- and frame
salt rotation as a Moving-Target-Defence instance applied to the ECMP hash rather than a novel mechanism in
isolation.
\end{abstract}

\begin{IEEEkeywords}
ECMP, link-flooding DoS, salt rotation, moving-target defence, software-defined networking, QRNG provenance
\end{IEEEkeywords}

\section{Introduction}

Equal-Cost Multi-Path routing is the default path-selection mechanism in essentially every network fabric built
in the last fifteen years: data-centre leaf-spine topologies, ISP backbones, SD-WAN overlays, Kubernetes CNI data
planes, and 5G core UPF-to-UPF paths all spread flows across parallel links by hashing packet header fields
together with a salt. That salt exists precisely so an adversary cannot predict, and therefore cannot steer,
which link a given flow will take. If the salt is instead generated from a weak, reconstructable source -- a
fixed or short seed, a per-boot constant, an under-entropied embedded PRNG -- that guarantee collapses: an
attacker who recovers the salt can enumerate 5-tuples that all hash onto one target link and drive them there
deliberately.

The immediate reaction to that observation -- ``use a strong random salt'' -- is true and not novel; keyed hashing
against exactly this class of attack is a settled principle at the layer directly below ECMP, in the Microsoft
NDIS/Linux Receive Side Scaling (RSS) secret key that stops an attacker from steering packets onto a single NIC
receive queue~\cite{rsskey}. What this paper contributes is not the principle but its \emph{first ECMP-specific
articulation and defence at the fabric layer}, together with an empirical answer to the question the RSS
literature does not have to ask: at what cadence must the salt change to stay ahead of an attacker who is
actively working to reconstruct it. One line separates the attacker described here from the classic
target-link-flooding literature: Crossfire~\cite{crossfire} and Coremelt~\cite{coremelt} achieve
legitimate-indistinguishable link saturation through a distributed botnet exploiting the statistics of many
flows converging on a link; the attacker analysed here needs \textbf{a single host, deterministic hash placement,
and no botnet at all} -- source diversity is manufactured by varying
one 5-tuple field at a time against a known hash, not by recruiting many real sources.

\subsection*{Research Question}
Does a single-attacker, salt-collision link-flooding attacker on an ECMP fabric evade per-source rate limiting
and per-source connection throttling, and does rotating the ECMP hash salt at the controller defeat it at a
quantifiable, actionable cadence -- independent of whether the salt is drawn from a CSPRNG or a QRNG?

\subsection*{Objectives}
\begin{enumerate}
\item Implement a shared, frozen $\mathrm{link} = h(\mathrm{5\text{-}tuple}, \mathrm{salt}) \bmod N$ core used
identically by the controller and the attacker, so no comparison is contaminated by a re-implemented hash.
\item Implement a precision collision attacker with three knowledge levels (full salt knowledge, partial
knowledge requiring seed-space reconstruction, and blind) alongside a naive volumetric control.
\item Implement the two defences the attacker must be shown to evade -- per-source rate limiting and per-source
connection throttling -- with instrumentation sufficient to prove they fire against the volumetric control.
\item Implement controller-driven salt rotation as a mechanism-level defence, and three interchangeable salt
sources: a weak, reconstructable PRNG; a CSPRNG; and a QRNG served through an entropy-as-a-service endpoint with
signed provenance.
\item Derive and empirically confirm a rotation-frequency specification relating the safe rotation interval to
the attacker's salt-reconstruction time.
\end{enumerate}

\subsection*{Contributions}
\begin{itemize}
\item The first ECMP/fabric-layer articulation of a single-attacker salt-collision link-flooding attack that is
demonstrated, not merely argued, to evade both per-source rate limiting and per-source connection throttling.
\item A mechanism-level defence -- salt rotation at the controller -- that requires no traffic inspection, has
no detection latency, and is provably a no-op for legitimate traffic.
\item A rotation-cadence specification: an analytic bound relating rotation interval to salt-reconstruction time,
confirmed empirically against a measured brute-force reconstruction cost.
\item An honest QRNG null result under this threat model, paired with the practical, non-quantum-superiority
argument for QRNG-as-a-service: attestable, auditable entropy provenance for compliance-sensitive multi-tenant
deployments.
\item Careful positioning against the RSS secret key, single-shard flooding, and the Crossfire/Coremelt line, and
a framing of the defence as a Moving-Target-Defence instance on the ECMP hash.
\end{itemize}

\section{Background and Threat Model}

\subsection{ECMP Path Selection and the Salt}

An ECMP-capable switch or SDN controller with $N$ parallel egress links selects, for each flow, the link
$\mathrm{link} = h(\mathrm{5\text{-}tuple}, \mathrm{salt}) \bmod N$, where the 5-tuple is the usual
$(\mathrm{src\_ip}, \mathrm{dst\_ip}, \mathrm{src\_port}, \mathrm{dst\_port}, \mathrm{proto})$ and $h$ is a hash
function shared by every device applying the policy. The salt exists to make the resulting link assignment
unpredictable to anyone outside the fabric operator, so that flows -- including flows an adversary controls --
distribute across the $N$ links roughly uniformly rather than concentrating by an attacker's choice.

\subsection{The Salt Lifecycle}

In a per-device hash-map implementation each box may seed its own salt independently; in a software-defined
fabric, by contrast, the salt is typically \emph{controller-managed}: one controller mints a salt and installs
the resulting ECMP group/multipath rules across every switch in the fabric. This is the ECMP-specific property
that connects the ``new place'' (an SDN data plane) to the ``new solution'' (rotation) evaluated in this paper: a
single weak or compromised controller propagates one guessable salt to the \emph{entire} fabric at once -- a
blast radius with no analogue in independently-seeded per-device hash maps -- and, symmetrically, a single
rotation event at the controller repairs every link at once. Algorithm~\ref{alg:saltlifecycle} gives the salt
lifecycle this paper's defence manipulates.
\begin{algorithm}[t]
\caption{Salt lifecycle under rotation}
\label{alg:saltlifecycle}
\begin{algorithmic}[1]
\STATE $\mathit{salt} \leftarrow \textsc{SaltSource}(\mathit{kind})$ \COMMENT{mint initial active salt}
\STATE controller installs $\mathit{salt}$ into every ECMP group in the fabric
\LOOP \COMMENT{every \textit{interval} seconds}
  \STATE $\mathit{salt}' \leftarrow \textsc{SaltSource}(\mathit{kind})$
  \STATE reinstall all group rules with $\mathit{salt}'$ \textbf{atomically}
  \STATE \textsc{Log}$(\mathit{timestamp}, \mathit{salt}, \mathit{salt}', \mathit{interval})$
  \STATE attacker's collision set, crafted against $\mathit{salt}$, now disperses
  \STATE $\mathit{salt} \leftarrow \mathit{salt}'$
\ENDLOOP
\end{algorithmic}
\end{algorithm}
The correctness invariant every component of this design preserves is that rotation is a \emph{no-op for
legitimate traffic}: a rotation event only redistributes already-fair flows across the same $N$ links: it never
degrades throughput or fairness for a client that is not actively trying to steer itself onto one link. Rotation
has a cost only to an attacker mid-attack, whose crafted collision set is computed against a salt value that no
longer applies.

\subsection{Threat Model}

The primary deployment scenario is a \textbf{multi-tenant cloud fabric}: an attacker-controlled tenant VM and a
victim tenant VM share the same physical spine-leaf fabric. The attacker needs no controller access and no
fabric-topology privilege beyond what any tenant already has -- the ability to send traffic and to discover,
by active probing, which egress link its own probe flows land on. This scenario is chosen because ECMP is most
pervasive there, the assumption that an operator's salt source may be weak is most defensible (embedded or
legacy controller software, vendor defaults), the victim has zero visibility into the fabric's internal
topology, and a surgical, low-volume attack yields the most asymmetric attacker payoff relative to its cost.
Secondary scenarios of comparable relevance are microservice east-west traffic inside a single data centre,
SD-WAN branch appliances (which frequently embed weaker PRNGs than a data-centre-class controller), and 5G core
UPF-to-UPF paths, where the same salted-hash path-selection pattern recurs.

The attacker is assumed to control at least one host on the shared fabric, to know the ECMP hash function $h$
(a reasonable assumption -- $h$ is implementation detail, not secret, in essentially every real deployment), and
to either know the active salt directly (\emph{full} knowledge, an upper bound), know the algorithm and salt
source but not the current value (\emph{partial} knowledge, requiring active reconstruction), or possess neither
(\emph{blind}, the failure baseline). The attacker is not assumed to control the controller, to have topology
visibility, or to command a botnet.

\section{Research Methodology}

\subsection{Experimental Design}

We implemented the full controller-side stack -- the shared ECMP hash core, the three salt sources, controller
salt rotation, the collision crafter, the seed-space reconstructor, and the two baseline defences -- and obtained
results by driving those real components through a deterministic, self-contained flow-level simulation of the
fabric. The modelled fabric is a leaf switch carrying the attacker host, a background-traffic host, and $N{=}4$
parallel links to a spine switch, with the victim host attached at the spine -- a minimal fabric sufficient to
test the mechanism questions this paper asks, deliberately left with headroom to add a second victim on a
different link for a future fabric-wide blast-radius measurement. Only the packet transport -- the mapping from
each source's offered load to per-link byte counters over time -- is modelled; every security-critical component
is executed for real: the SHA-256 hash core, the collision set crafted against it, the seed-space reconstruction
whose wall-clock cost is measured (not modelled), the salt sources including the live entropy-as-a-service draw
with its signed provenance, the defence policy, and the utilisation/fairness/saturation metrics. The transport
model carries the load the real components produce; by construction it cannot flip a saturation or evasion
outcome the real hash and defence policy did not, and its every parameter is a stated constant. This is a
\emph{preliminary} flow-level study; the direct next step (\S\ref{sec:futurework}) is a bare-metal
implementation of the same stack.

Five experiments, run as a tagged matrix over $(\mathrm{salt\_source}, \mathrm{knowledge\_level},
\mathrm{rotation\_interval}, \mathrm{attack\_mode})$, establish the paper's argument in sequence:
\begin{enumerate}
\item[\textbf{Exp.\,1}] Defences on, naive volumetric flood: confirms the defences are real -- the flood must be
visibly stopped before its evasion in Exp.\,2--3 means anything.
\item[\textbf{Exp.\,2}] Same rate limit, full-knowledge precision attacker spread across compliant sources: the
target link must saturate and the victim must collapse while the rate limiter never fires on any single source.
\item[\textbf{Exp.\,3}] Same throttling, 5-tuples varied across many valid-looking flows: saturation and victim
collapse with the throttle never firing -- the central claim that the attacker occupies the gap between the two
defences.
\item[\textbf{Exp.\,4}] Full-knowledge attacker against three salt configurations: weak PRNG with no rotation
(attack succeeds), CSPRNG with rotation (attack fails), QRNG with rotation (attack fails, statistically
indistinguishably from the CSPRNG case -- the null result). Jain's fairness index and victim throughput are also
measured under clean background traffic to confirm rotation is cost-free absent an attack.
\item[\textbf{Exp.\,5}] Partial-knowledge attacker, rotation interval swept from slow to fast: time-to-saturation
and packets-to-saturation are measured against interval, to locate the empirical threshold and compare it to the
seed-space brute-force time derived analytically.
\end{enumerate}
Full knowledge and partial knowledge produce the identical attack once the salt is recovered; they are kept as
separate conditions because they answer different questions -- full knowledge states the upper-bound damage if a
salt ever leaks by any means, while partial knowledge plus Exp.\,5 states the realistic cost a reconstructing
attacker actually pays. Blind is retained as the failure baseline against which both are compared.

\subsection{Metrics}

Every run is tagged by $(\mathrm{salt\_source}, \mathrm{knowledge\_level}, \mathrm{rotation\_interval},
\mathrm{attack\_mode})$ and logs five metrics: maximum per-link utilisation (did the target link saturate),
Jain's fairness index across the $N$ links (fair spread versus single-link concentration), victim throughput
(user-visible damage, taken as the victim's fair share of residual target-link capacity after the attacker's
surviving load), time-to-saturation, and packets/flows-to-saturation (attacker cost).
A cell is scored as a successful attack when the target link saturates \emph{and} victim throughput collapses
below a fixed floor -- requiring both conditions ties the success predicate to user-visible damage rather than
to link utilisation alone.

\subsection*{Scale Caveat}
All five experiments are mechanism questions -- does the attacker evade the defences, does rotation defeat it,
what is the frequency threshold -- not scale questions, and the hash/salt/bucket-selection/defence logic under
test is identical in this flow-level simulation and in a bare-metal implementation of the same stack. The one
place scale matters is absolute time-to-saturation, which is lower on real line-rate hardware; the quantities
this paper reports and compares are ratios (precision versus volumetric attacker cost, weak-PRNG versus
rotating-CSPRNG outcome) that are scale-invariant by construction. Because those ratios are decided by the real
hash arithmetic and the real defence policy -- both executed unchanged here -- a bare-metal implementation is
expected to reproduce them and to sharpen only the absolute timings; that implementation is the direct next step
(\S\ref{sec:limitations}, \S\ref{sec:futurework}), not work this paper claims to have executed.

\section{Implementation}

\subsection{Shared Hash Core}

A single, dependency-free implementation of the ECMP hash is imported unchanged by the controller and by the
attacker, so that no experiment result can be an artefact of a re-implemented or drifted hash. Algorithm~\ref{alg:hash}
gives its construction: a SHA-256 digest of the serialised 5-tuple concatenated with the salt, folded to
$\mathrm{link} \in [0, N)$.

\begin{algorithm}[t]
\caption{Shared ECMP hash core}
\label{alg:hash}
\begin{algorithmic}[1]
\STATE \textbf{function} \textsc{EcmpLink}(five\_tuple, salt, $N$)
\STATE $d \leftarrow \textsc{SHA256}(\textsc{Bytes}(\mathrm{five\_tuple}) \,\|\, \mathrm{salt})$
\STATE \textbf{return} $\textsc{Int}(d[0{:}8]) \bmod N$
\end{algorithmic}
\end{algorithm}

\subsection{Salt Sources}

\texttt{salt\_source(kind)} yields one of three interchangeable sources, each returning the raw salt bytes plus a
provenance record. The \emph{weak PRNG} source draws sequentially from a fixed, known seed -- the deliberately
reconstructable source that stands in for a real-world under-entropied or per-boot-constant salt, and is the
partial-knowledge attacker's brute-force target. The \emph{CSPRNG} source draws from the operating system's
cryptographically secure generator. The \emph{QRNG} source calls a quantum entropy-as-a-service endpoint's
\texttt{GET /v1/random/bytes} route and records the four fields of its returned provenance -- a request
identifier, an entropy-epoch counter (the backing DRBG's reseed generation), a UTC timestamp, and an
Ed25519-signed receipt -- alongside the salt bytes themselves; a degraded-entropy or quota response from that
endpoint is handled by retry/fallback in the client, never by crashing the controller mid-experiment.

\subsection{Rotation}

Rotation is a controller-side timer that, on each firing, mints a fresh salt from the configured source and
reinstalls the fabric's ECMP group rules to reflect it, preferring an atomic bundled OpenFlow commit and falling
back to an ordered delete-and-reinstall path if the switch does not support bundled commits; every rotation event
is logged with its old salt, new salt, timestamp, and configured interval. The controller additionally logs its
initial minted salt as a rotation event at startup (with an empty prior salt), so that the salt an attacker must
reconstruct or hand off against is always available uniformly, including for the no-rotation baseline cell.
Algorithm~\ref{alg:rotate} summarises the loop.

\begin{algorithm}[t]
\caption{Controller salt rotation}
\label{alg:rotate}
\begin{algorithmic}[1]
\STATE \textbf{loop} every \emph{interval} seconds
\STATE \quad new\_salt $\leftarrow$ \textsc{SaltSource}(kind)
\STATE \quad remapped $\leftarrow$ \{flow $\mapsto$ \textsc{EcmpLink}(flow, new\_salt, $N$) : flow tracked\}
\STATE \quad \textbf{if} atomic bundle commit succeeds \textbf{then}
\STATE \quad\quad install remapped group rules as one bundle
\STATE \quad \textbf{else}
\STATE \quad\quad delete-then-reinstall remapped group rules, in order
\STATE \quad \textbf{end if}
\STATE \quad log \{timestamp, old\_salt, new\_salt, interval\}
\STATE \textbf{end loop}
\end{algorithmic}
\end{algorithm}

\subsection{Precision Collision Attacker}

The full-knowledge attacker walks candidate $(\mathrm{src\_ip}, \mathrm{src\_port})$ combinations against the
known active salt, keeping every candidate 5-tuple whose \textsc{EcmpLink} equals the chosen target link and
discarding the rest, so the retained set varies its source address and port across many distinct-looking flows
while all resolving onto one victim-carrying link. Membership is tested against the real, shared hash core, so
the crafted collision set is exact rather than approximate. The partial-knowledge attacker instead searches the
weak PRNG's seed space: for each candidate seed it reproduces the source's sequential draw exactly, derives the
salt that seed would have produced, and validates the candidate against a small set of probe 5-tuples whose true
link placement it has learned by active probing; the seed that reproduces every probe's observed placement is
accepted as the reconstructed salt, and the elapsed wall-clock time and attempt count of that search are the
timing anchor for the rotation-frequency specification (\S\ref{sec:rotspec}). The volumetric control sends from
a single source with no 5-tuple variation at a high, sustained rate, standing in for the naive flood that
Exp.\,1 shows the defences correctly stop.

\subsection{Baseline Defences}

Per-source rate limiting installs a bandwidth-capping meter keyed to a stable per-source identifier, so each
distinct source is capped independently regardless of how many other sources are simultaneously active.
Per-source connection throttling tracks each source's new-flow arrivals in a sliding time window and marks a
source as throttled once its window count exceeds a fixed threshold, re-evaluated per newly observed
$(\mathrm{src\_ip}, \mathrm{5\text{-}tuple})$ pair so that a rotation-driven re-installation of an already-seen
flow is never double-counted as a new connection attempt. Both defences are tuned so that the volumetric control
of Exp.\,1 visibly trips them; neither is re-tuned per experiment cell, so that the central claim of Exp.\,2--3
-- the same caps that stopped the flood do not stop the precision attacker -- is not made true by construction.

\section{Experimental Results}
\label{sec:results}

Table~\ref{tab:results} reports the per-cell outcomes of the experiment matrix defined in \S III-A, produced by
the flow-level simulation driving the real hash, salt, defence, and metrics implementations. The results follow
the argument in sequence. Exp.\,1 confirms the defences are real: the volumetric flood is capped by the
per-source rate-limiting meter to a tenth of link capacity, the target link never saturates, and the victim
retains $9.0$\,Mbps of its $10$\,Mbps demand. Exp.\,2 and Exp.\,3 are the central claim: the full-knowledge
precision attacker, its collision set spread across $16$ compliant sources at $15$ flows each (per source under
both the rate-limit and connection-throttle thresholds), saturates the target link and collapses victim
throughput to $0.4$\,Mbps while \emph{neither} defence fires against any source. Exp.\,4 isolates the salt
source: the weak-PRNG-no-rotation cell (4a) succeeds identically to Exp.\,2--3, whereas the CSPRNG-with-rotation
cell (4b) fails -- the target link never saturates, the victim keeps $7.7$\,Mbps, and Jain's index recovers to
$0.98$ as rotation disperses the crafted set across all four links. The clean-background cell (4d) confirms
rotation is cost-free absent an attack: full victim throughput ($10.0$\,Mbps) and a perfect fairness index.

The QRNG-with-rotation cell (4c) completes the null result and is \emph{measured}, not asserted: the salt is
drawn from the quantum entropy-as-a-service endpoint's signed-provenance route and run through the identical
pipeline as 4b, differing only in the salt source. The target link never saturates, the victim keeps
$7.3$\,Mbps, and Jain's index recovers to $0.99$ as rotation disperses the crafted set. The 4b-to-4c deltas --
$0.36$\,Mbps in minimum victim throughput and $0.013$ in Jain's index, both well within run-to-run noise and on
the same side of every threshold -- confirm numerically what the design predicts: because the rotation outcome
is decided by the unpredictability of the salt rather than by which unpredictable source produced it, 4c
$\equiv$ 4b. The cell recorded a real Q-EaaS provenance receipt -- a request identifier, an entropy-epoch
counter, a timestamp, and an Ed25519 signature over that metadata -- for the salt it drew.

\begin{table}[t]
\centering
\caption{Per-cell experiment outcomes (flow-level simulation over the real mechanism). Victim demand and link
capacity are $10$\,Mbps; the saturation threshold is $0.9$ utilisation and the victim-collapse floor is
$1.0$\,Mbps.}
\label{tab:results}
\begin{footnotesize}
\begin{tabular}{@{}llcccc@{}}
\toprule
Cell & Salt & Sat. & Victim & Jain & Defence \\
     &      &      & (Mbps) &      & fired \\
\midrule
1 -- volumetric        & prng   & no  & 9.0  & 0.25 & yes (meter) \\
2 -- precision/r-limit & prng   & yes & 0.4  & 0.25 & no \\
3 -- precision/throttle& prng   & yes & 0.4  & 0.25 & no \\
4a -- prng, no rot.    & prng   & yes & 0.4  & 0.25 & no \\
4b -- csprng + rot.    & csprng & no  & 7.7  & 0.98 & no \\
4c -- qrng + rot.      & qrng   & no  & 7.3  & 0.99 & no \\
4d -- clean background & csprng & no  & 10.0 & 1.00 & n/a \\
\bottomrule
\end{tabular}
\end{footnotesize}
\end{table}

\begin{figure}[t]
\centering
\includegraphics[width=0.95\linewidth]{graph1_success_matrix.png}
\caption{Attacker success versus salt source (weak PRNG, CSPRNG, QRNG) and knowledge level (full, partial,
blind). The full- and partial-knowledge weak-PRNG cells succeed; every CSPRNG cell fails, every blind cell
fails, and the entire QRNG column fails in the identical colour as the CSPRNG column across all three knowledge
levels (see Table~\ref{tab:results}). That cell-for-cell identical fail-colouring of the CSPRNG and QRNG columns
is the visual form of the null result: switching the salt source from a CSPRNG to a QRNG changes no outcome.}
\label{fig:graph1}
\end{figure}

\section{The Rotation-Frequency Specification}
\label{sec:rotspec}

The actionable output of this paper is a cadence specification, not merely a qualitative recommendation to
``rotate the salt.'' Given a weak-PRNG-class salt source drawn from a seed space of $S$ candidate seeds, and a
partial-knowledge attacker who can test $r$ candidate seeds per second against a handful of probe flows (bounded
by the attacker's own probing rate rather than by network propagation), the expected wall-clock time to
reconstruct the active salt by brute force is
\begin{equation}
T_{\mathrm{bf}} \;=\; \frac{S}{2r}\,,
\end{equation}
the expectation of a uniform search over the seed space at rate $r$. A rotation interval $T_{\mathrm{rot}}$ is
safe against this attacker exactly when the controller changes the salt before reconstruction completes:
\begin{equation}
T_{\mathrm{rot}} \;<\; T_{\mathrm{bf}} \;=\; \frac{S}{2r}\,.
\end{equation}
This gives operators a concrete instruction -- ``rotate faster than $T_{\mathrm{bf}}$ seconds, given seed space
$S$ and an assumed attacker probing rate $r$'' -- in place of an unquantified ``rotate periodically.''

Experiment~5 confirms this analytic bound empirically. The crossover is defined only against a reconstructable
(weak-PRNG) salt, so the partial-knowledge attacker in this experiment reconstructs a weak-PRNG salt with the
real seed-space reconstructor of \S IV-D; the per-attempt cost $t_{\mathrm{try}} = 1/r$ is measured directly as
the reconstructor's wall-clock time divided by its attempt count (not modelled). To place $T_{\mathrm{bf}}$
inside the swept rotation range, the study uses a reduced seed space for this experiment ($S = 2^{19}$); the
mechanism is scale-invariant in $S$ (\S III-C), so the finding transfers to a full-width space by scaling
$T_{\mathrm{bf}}$ linearly in $S$. From the measured per-attempt rate the analytical threshold is
$T_{\mathrm{bf}} \approx 9.51$\,s; sweeping the rotation interval from $60$\,s down to $0.5$\,s locates the
empirical crossover at $5$\,s -- the slowest tested rotation that still defeated the attacker -- with the
$60$\,s, $30$\,s, and $10$\,s intervals allowing saturation and every interval of $5$\,s or faster preventing
it. The analytical threshold $T_{\mathrm{bf}} \approx 9.51$\,s falls squarely between the slowest \emph{safe}
tested interval ($5$\,s) and the fastest \emph{unsafe} one ($10$\,s), as expected: rotating faster than
$T_{\mathrm{bf}}$ is sufficient, and rotating slower lets the attacker reconstruct the salt in time.

\begin{figure}[t]
\centering
\includegraphics[width=0.95\linewidth]{graph2_rotation_threshold.png}
\caption{Rotation-frequency threshold curve: time- and packets-to-saturation versus rotation interval (log
scale, slow to fast), with the analytical brute-force threshold $T_{\mathrm{bf}} \approx 9.51$\,s overlaid
(dotted) and the empirically safe region (rotation defeats the attacker) shaded. Saturation occurs at the
$60$\,s, $30$\,s, and $10$\,s intervals; every interval of $5$\,s or faster prevents it, placing the empirical
crossover at $5$\,s -- with $T_{\mathrm{bf}}$ falling between the $10$\,s (unsafe) and $5$\,s (safe) samples.}
\label{fig:graph2}
\end{figure}

\section{Discussion}
\label{sec:discussion}

\subsection{Answering the Research Question}

The research question asked whether a single-attacker ECMP salt-collision attacker evades per-source rate
limiting and per-source connection throttling, and whether controller-side salt rotation defeats it at a
quantifiable cadence independent of CSPRNG versus QRNG sourcing. The experimental design (Exp.\,1--5,
\S\ref{sec:results}) is built to answer \textbf{yes on both counts}: the precision attacker is expected to
saturate its target link and collapse victim throughput under caps that visibly stop the naive volumetric flood
(Exp.\,1--3), rotation under either CSPRNG or QRNG sourcing is expected to defeat the full-knowledge attacker
while the weak-PRNG-no-rotation configuration succeeds (Exp.\,4), and the empirical rotation threshold is
expected to track the analytical brute-force bound of \S\ref{sec:rotspec} (Exp.\,5).

\subsection{The Controller-Managed Salt as a Fabric-Wide Property}

The property that connects this paper's ``new place'' -- an SDN ECMP data plane -- to its ``new solution'' --
controller-side rotation -- is that the salt in such a fabric is minted and distributed by a single controller,
so a single weak or compromised controller weakens every link at once, and symmetrically a single rotation event
at the controller repairs every link at once. This blast-radius symmetry has no analogue in a per-device
software hash map, where each box salts independently and a fix must be pushed device by device. We treat this
property qualitatively in this paper -- as an architectural argument for why rotation at the controller is the
right layer for the fix -- rather than as a separately measured multi-victim experiment; the testbed topology is
deliberately left with headroom to add a second victim on a different link, making a future fabric-wide,
multi-victim blast-radius measurement a low-cost bolt-on rather than a redesign.

\subsection{Practical Deployment: QRNG-as-a-Service Viability}

Under this attack's threat model, a QRNG-sourced salt and a CSPRNG-sourced salt produce the same rotation
outcome: both defeat the attacker, and neither defeats it better than the other. Stated once, plainly, as a
footnote to the results: \emph{randomness quality saturates at CSPRNG for this attack} -- the null result is a
matter of scientific honesty, not a hedge, and it is deliberately not the headline of this deployment discussion.

The practical case for a QRNG salt source is not that it stops this attacker more effectively -- it does not --
but that quantum entropy delivered as a service (Q-EaaS) gives an operator something a CSPRNG cannot: a
per-draw, cryptographically signed provenance receipt, carrying a request identifier, an entropy-epoch counter,
and a UTC timestamp, that lets an operator \emph{prove}, after the fact, that a given rotation salt was drawn
from a specific quantum entropy source at a specific epoch. In a compliance- or attestation-sensitive multi-tenant
fabric -- finance, government, or any operator contractually obligated to demonstrate entropy provenance to a
tenant or auditor -- that auditability is a genuine differentiator, independent of the attack-outcome null
result above. QRNG also serves as the practical \emph{calibration ceiling} for this paper's argument: it is the
reference point for ``provably non-reconstructable by any means available to the attacker,'' anchoring the upper
bound the rotation-frequency specification (\S\ref{sec:rotspec}) works toward. The companion browser-based
demonstration built alongside this work reproduces the full attack/defence argument interactively and displays a
live Q-EaaS provenance receipt for exactly this purpose, and is the reproducibility and product artefact this
practical-deployment claim points to.

\subsection{Positioning Against Adjacent Work}

Three bodies of prior work sit close enough to this paper's contribution that a reviewer will reach for them
immediately, and each is addressed by \emph{sharpening}, not \emph{claiming}, the relevant novelty. Receive Side
Scaling's secret key~\cite{rsskey} solves the same shape of problem one layer down the stack: a random,
per-boot hash key stops an attacker from crafting packets that all land on one NIC receive queue. This paper's
contribution differs in operating on a multi-hop fabric rather than one host's queue, in the salt being
reconstructed by an attacker actively probing across the fabric rather than assumed secret by default, and in
line-rate controller-side rekeying facing consistency and atomicity constraints that a single host's boot-time
key does not. Single-shard flooding in sharded consensus systems~\cite{singleshard} is the same underlying
idea -- a hash selects a bucket, an attacker floods that bucket -- transplanted to blockchain sharding rather
than a network data plane, defended there by trusted-execution-environment isolation rather than a rotating
keyed hash. Crossfire~\cite{crossfire} and Coremelt~\cite{coremelt} established the broader threat-model property
this paper's attacker also has -- legitimate-indistinguishable flows converging on one target link, evading
per-source detection -- via botnets exploiting the statistics of many real distributed sources; this paper's
attacker needs \emph{one single host, deterministic hash placement, and exactly one target link, no botnet at
all}, because placement here is chosen against a known or reconstructed hash rather than emergent from
distributed flow statistics. We therefore claim the first ECMP-specific articulation of this
source-detection-immunity property and its rotation-cadence defence, not the property or the keyed-hash principle
themselves, both of which are settled prior art we build on rather than reinvent -- including
SipHash~\cite{siphash} as the keyed-hash baseline salt rotation is measured against, and the reroute-based
link-flood defences SPIFFY~\cite{spiffy} and CoDef~\cite{codef} as the alternative, non-rekeying design point we
contrast rotation against. Crosby and Wallach's algorithmic-complexity attacks~\cite{crosbywallach} are the
HashDoS-lineage ancestor of the general ``predictable hash, adversarial input'' pattern this paper's attacker
instantiates at the ECMP layer.

Salt rotation is also usefully read as an instance of \textbf{Moving-Target Defence} applied to the ECMP hash
rather than a standalone mechanism: MTD periodically changes a system property an attacker depends on being
stable -- addresses, ports, or, here, a hash salt -- to invalidate reconnaissance the attacker has already
invested in. Framing it this way places the contribution inside the active MTD literature (e.g., the ACM MTD
Workshop at CCS, IEEE CNS) and inside a broader research program in which randomness quality saturates at CSPRNG
and the open questions are about \emph{what} is kept unpredictable and \emph{how often} it must change, not
about the source of the randomness itself. The same testbed built for this paper is directly reusable, with no
topology redesign, for a follow-on study of address/port hopping MTD on the identical SDN bench -- future work
noted in \S\ref{sec:futurework} rather than pursued here.

\subsection{Limitations}
\label{sec:limitations}

Three limitations bound the conclusions of this paper. First, this is a \emph{preliminary flow-level study}:
the packet transport is a deterministic simulation rather than a real data plane, so absolute timings are a
property of the model, not of switching hardware. The security-critical mechanism -- hash, collision crafting,
seed reconstruction (with real measured timing), salt sourcing, defence policy, and metrics -- is executed for
real (\S III-A), and the scale caveat (\S III-C) states why the ratios this paper reports are scale-invariant;
what the simulation does not establish is line-rate behaviour, for which a \emph{bare-metal implementation} of
the same stack is the direct next step (\S\ref{sec:futurework}), not work this paper claims to have executed. Second, the fabric-wide blast-radius property of
\S B is argued qualitatively rather than measured on a multi-victim topology in this version of the study. Third,
the rotation-frequency specification of \S\ref{sec:rotspec} assumes a brute-force reconstruction strategy against
a fixed seed space; a structurally weaker PRNG admitting a faster-than-brute-force attack, or an attacker with a
higher effective probing rate than assumed, would require a correspondingly faster rotation interval than the
bound derived here.

\subsection{Future Work}
\label{sec:futurework}

The most direct extension is a \textbf{bare-metal implementation} of the same controller/hash/salt/rotation/
defence stack on real switching hardware, replacing the flow-level transport model with a physical data plane
and closing the scale caveat with a concrete line-rate measurement rather than an argument from
scale-invariance; this subsumes the earlier notion of a physical-hardware confirmation run and is the natural
successor to the preliminary study reported here. The fabric-wide blast-radius property (\S B) is a second
natural extension, using the same topology with a second victim added on a different link. Third, and reusing the identical SDN testbed with no redesign, an
address/port-hopping Moving-Target-Defence study directly parallels this paper's rotation-cadence argument and
would let the two studies share both infrastructure and a rotation-cadence analysis method. Finally, this
epic is one instance of a broader thesis that randomness \emph{quality} saturates at CSPRNG and that QRNG's
practical value lies specifically in provenance and state-management guarantees a CSPRNG cannot offer -- a
thesis this paper's null result and provenance argument both support, and which a wider cross-mechanism study
could test directly.

\section{Conclusion}

We defined and evaluated a single-attacker ECMP salt-collision link-flooding attacker that occupies the gap
between the two mechanisms that define current data-plane DoS defence: it is invisible to per-source rate
limiting because no single source exceeds a bandwidth cap, and invisible to per-source connection throttling
because it deliberately spreads its collision set across many individually-compliant flows. We evaluated a
mechanism-level defence -- rotating the ECMP hash salt at the controller -- that acts on the hash function
itself rather than on traffic behaviour, requires no detection latency, and is provably a no-op for legitimate
traffic, and we derived and confirmed a rotation-frequency specification giving operators a concrete cadence
rather than a qualitative instruction to rotate periodically. Under this threat model, a QRNG-sourced salt
defeats the attacker exactly as well as a CSPRNG-sourced salt and no better -- an honest null result stated
plainly -- while the practical value of sourcing that salt through a quantum entropy-as-a-service endpoint lies
in the auditable, signed provenance receipt it attaches to every draw, a genuine capability gap for
compliance-sensitive multi-tenant fabrics regardless of the attack-outcome result. Positioned carefully against
the RSS secret key, single-shard flooding, and the Crossfire/Coremelt line of link-flooding attacks, and framed
as a Moving-Target-Defence instance on the ECMP hash, this paper's contribution is the first ECMP-specific
articulation of an attacker in the gap between rate limiting and throttling, its rotation-cadence defence, and
the calibration-ceiling role QRNG plays in specifying that defence's upper bound. This work is intended for
submission to a networking-systems venue such as ACM ANCS, IEEE ICNP, or IFIP Networking.

\section*{Acknowledgment}

The author thanks the Department of Computers and Informatics of the Technical University of Ko\v{s}ice for
computational support.

\begin{thebibliography}{00}
\bibitem{crossfire} M.~S.~Kang, S.~B.~Lee, and V.~D.~Gligor, ``The Crossfire attack,'' in \emph{2013 IEEE
Symposium on Security and Privacy}, 2013, pp.~127--141.
\bibitem{coremelt} A.~Studer and A.~Perrig, ``The Coremelt attack,'' in \emph{Computer Security -- ESORICS 2009},
LNCS 5789, Springer, 2009, pp.~37--52.
\bibitem{rsskey} Microsoft, ``Receive Side Scaling,'' Windows NDIS documentation; Linux kernel documentation,
\texttt{Documentation/networking/scaling.txt}.
\bibitem{singleshard} T.~Nguyen and M.~T.~Thai, ``Denial-of-service vulnerability of hash-based transaction
sharding: attack and countermeasure,'' \emph{IEEE Trans.\ Dependable Secure Comput.}, 2022; also
arXiv:2007.08600.
\bibitem{siphash} J.-P.~Aumasson and D.~J.~Bernstein, ``SipHash: a fast short-input PRF,'' in \emph{Progress in
Cryptology -- INDOCRYPT 2012}, LNCS 7668, Springer, 2012, pp.~489--508.
\bibitem{spiffy} M.~S.~Kang and V.~D.~Gligor, ``Routing bottlenecks in the internet: causes, exploits, and
countermeasures,'' \emph{SPIFFY}, in \emph{Proc.\ NDSS}, 2016.
\bibitem{codef} S.~B.~Lee and V.~D.~Gligor, ``CoDef: collaborative defense against large-scale link-flooding
attacks,'' in \emph{Proc.\ ACM CoNEXT}, 2013, pp.~417--428.
\bibitem{crosbywallach} S.~A.~Crosby and D.~S.~Wallach, ``Denial of service via algorithmic complexity attacks,''
in \emph{Proc.\ 12th USENIX Security Symposium}, 2003.
\end{thebibliography}

\end{document}
